% Options for packages loaded elsewhere
\PassOptionsToPackage{unicode}{hyperref}
\PassOptionsToPackage{hyphens}{url}
%
\documentclass[
]{article}
\usepackage{amsmath,amssymb}
\usepackage{iftex}
\ifPDFTeX
  \usepackage[T1]{fontenc}
  \usepackage[utf8]{inputenc}
  \usepackage{textcomp} % provide euro and other symbols
\else % if luatex or xetex
  \usepackage{unicode-math} % this also loads fontspec
  \defaultfontfeatures{Scale=MatchLowercase}
  \defaultfontfeatures[\rmfamily]{Ligatures=TeX,Scale=1}
\fi
\usepackage{lmodern}
\ifPDFTeX\else
  % xetex/luatex font selection
\fi
% Use upquote if available, for straight quotes in verbatim environments
\IfFileExists{upquote.sty}{\usepackage{upquote}}{}
\IfFileExists{microtype.sty}{% use microtype if available
  \usepackage[]{microtype}
  \UseMicrotypeSet[protrusion]{basicmath} % disable protrusion for tt fonts
}{}
\makeatletter
\@ifundefined{KOMAClassName}{% if non-KOMA class
  \IfFileExists{parskip.sty}{%
    \usepackage{parskip}
  }{% else
    \setlength{\parindent}{0pt}
    \setlength{\parskip}{6pt plus 2pt minus 1pt}}
}{% if KOMA class
  \KOMAoptions{parskip=half}}
\makeatother
\usepackage{xcolor}
\usepackage{color}
\usepackage{fancyvrb}
\newcommand{\VerbBar}{|}
\newcommand{\VERB}{\Verb[commandchars=\\\{\}]}
\DefineVerbatimEnvironment{Highlighting}{Verbatim}{commandchars=\\\{\}}
% Add ',fontsize=\small' for more characters per line
\newenvironment{Shaded}{}{}
\newcommand{\AlertTok}[1]{\textcolor[rgb]{1.00,0.00,0.00}{\textbf{#1}}}
\newcommand{\AnnotationTok}[1]{\textcolor[rgb]{0.38,0.63,0.69}{\textbf{\textit{#1}}}}
\newcommand{\AttributeTok}[1]{\textcolor[rgb]{0.49,0.56,0.16}{#1}}
\newcommand{\BaseNTok}[1]{\textcolor[rgb]{0.25,0.63,0.44}{#1}}
\newcommand{\BuiltInTok}[1]{\textcolor[rgb]{0.00,0.50,0.00}{#1}}
\newcommand{\CharTok}[1]{\textcolor[rgb]{0.25,0.44,0.63}{#1}}
\newcommand{\CommentTok}[1]{\textcolor[rgb]{0.38,0.63,0.69}{\textit{#1}}}
\newcommand{\CommentVarTok}[1]{\textcolor[rgb]{0.38,0.63,0.69}{\textbf{\textit{#1}}}}
\newcommand{\ConstantTok}[1]{\textcolor[rgb]{0.53,0.00,0.00}{#1}}
\newcommand{\ControlFlowTok}[1]{\textcolor[rgb]{0.00,0.44,0.13}{\textbf{#1}}}
\newcommand{\DataTypeTok}[1]{\textcolor[rgb]{0.56,0.13,0.00}{#1}}
\newcommand{\DecValTok}[1]{\textcolor[rgb]{0.25,0.63,0.44}{#1}}
\newcommand{\DocumentationTok}[1]{\textcolor[rgb]{0.73,0.13,0.13}{\textit{#1}}}
\newcommand{\ErrorTok}[1]{\textcolor[rgb]{1.00,0.00,0.00}{\textbf{#1}}}
\newcommand{\ExtensionTok}[1]{#1}
\newcommand{\FloatTok}[1]{\textcolor[rgb]{0.25,0.63,0.44}{#1}}
\newcommand{\FunctionTok}[1]{\textcolor[rgb]{0.02,0.16,0.49}{#1}}
\newcommand{\ImportTok}[1]{\textcolor[rgb]{0.00,0.50,0.00}{\textbf{#1}}}
\newcommand{\InformationTok}[1]{\textcolor[rgb]{0.38,0.63,0.69}{\textbf{\textit{#1}}}}
\newcommand{\KeywordTok}[1]{\textcolor[rgb]{0.00,0.44,0.13}{\textbf{#1}}}
\newcommand{\NormalTok}[1]{#1}
\newcommand{\OperatorTok}[1]{\textcolor[rgb]{0.40,0.40,0.40}{#1}}
\newcommand{\OtherTok}[1]{\textcolor[rgb]{0.00,0.44,0.13}{#1}}
\newcommand{\PreprocessorTok}[1]{\textcolor[rgb]{0.74,0.48,0.00}{#1}}
\newcommand{\RegionMarkerTok}[1]{#1}
\newcommand{\SpecialCharTok}[1]{\textcolor[rgb]{0.25,0.44,0.63}{#1}}
\newcommand{\SpecialStringTok}[1]{\textcolor[rgb]{0.73,0.40,0.53}{#1}}
\newcommand{\StringTok}[1]{\textcolor[rgb]{0.25,0.44,0.63}{#1}}
\newcommand{\VariableTok}[1]{\textcolor[rgb]{0.10,0.09,0.49}{#1}}
\newcommand{\VerbatimStringTok}[1]{\textcolor[rgb]{0.25,0.44,0.63}{#1}}
\newcommand{\WarningTok}[1]{\textcolor[rgb]{0.38,0.63,0.69}{\textbf{\textit{#1}}}}
\usepackage{longtable,booktabs,array}
\usepackage{calc} % for calculating minipage widths
% Correct order of tables after \paragraph or \subparagraph
\usepackage{etoolbox}
\makeatletter
\patchcmd\longtable{\par}{\if@noskipsec\mbox{}\fi\par}{}{}
\makeatother
% Allow footnotes in longtable head/foot
\IfFileExists{footnotehyper.sty}{\usepackage{footnotehyper}}{\usepackage{footnote}}
\makesavenoteenv{longtable}
\setlength{\emergencystretch}{3em} % prevent overfull lines
\providecommand{\tightlist}{%
  \setlength{\itemsep}{0pt}\setlength{\parskip}{0pt}}
\setcounter{secnumdepth}{-\maxdimen} % remove section numbering
\ifLuaTeX
  \usepackage{selnolig}  % disable illegal ligatures
\fi
\IfFileExists{bookmark.sty}{\usepackage{bookmark}}{\usepackage{hyperref}}
\IfFileExists{xurl.sty}{\usepackage{xurl}}{} % add URL line breaks if available
\urlstyle{same}
\hypersetup{
  hidelinks,
  pdfcreator={LaTeX via pandoc}}

\author{}
\date{}

\begin{document}

\hypertarget{bip-353-dns-payment-instructions}{%
\section{BIP-353 DNS Payment
Instructions}\label{bip-353-dns-payment-instructions}}

SatsPath can resolve \textbf{DNSSEC-backed BIP-353 payment instructions}
for domains the receiver controls. A receiver publishes a
\texttt{bitcoin:} URI in DNS; any SatsPath-compatible client can then
resolve \texttt{₿user@domain} without the sender having the receiver
locally registered.

\begin{quote}
\textbf{Mainnet Preview only.} This is a resolver/publisher layer.
SatsPath does not execute mainnet payments. It resolves, verifies,
routes, and displays payment instructions --- it never pays, signs, or
broadcasts.
\end{quote}

\hypertarget{what-bip-353-is}{%
\subsection{What BIP-353 is}\label{what-bip-353-is}}

\href{https://github.com/bitcoin/bips/blob/master/bip-0353.mediawiki}{BIP-353}
defines human-readable Bitcoin payment names of the form
\texttt{₿user@domain}, resolved through DNSSEC-signed DNS \texttt{TXT}
records.

\begin{itemize}
\tightlist
\item
  Records live at:
  \texttt{\textless{}user\textgreater{}.user.\_bitcoin-payment.\textless{}domain\textgreater{}}
\item
  The \texttt{TXT} RDATA reconstructs into a \texttt{bitcoin:} URI
  (BIP-321 payment instructions).
\item
  All payment instructions \textbf{must} be DNSSEC-signed.
\end{itemize}

Example:

\begin{verbatim}
rodrigo.user._bitcoin-payment.satspath.dev TXT "bitcoin:?lno=lno1..."
\end{verbatim}

\hypertarget{how-satspath-uses-it}{%
\subsection{How SatsPath uses it}\label{how-satspath-uses-it}}

\begin{verbatim}
₿rodrigo@satspath.dev  ──parse──▶  rodrigo.user._bitcoin-payment.satspath.dev
                                          │ DNSSEC-validated TXT lookup
                                          ▼
                          bitcoin:?lno=lno1...   (BIP-321 URI)
                                          │ parse + screen
                                          ▼
                      Bip353Resolution → Mainnet-Preview QuoteResponse
\end{verbatim}

CLI:

\begin{Shaded}
\begin{Highlighting}[]
\ExtensionTok{satspath}\NormalTok{ dns resolve ₿rodrigo@satspath.dev}
\ExtensionTok{satspath}\NormalTok{ dns resolve rodrigo@satspath.dev }\AttributeTok{{-}{-}json}
\end{Highlighting}
\end{Shaded}

Resolution rules SatsPath enforces (per BIP-353):

\begin{itemize}
\tightlist
\item
  Reconstruct a \texttt{TXT} record by concatenating its
  \texttt{\textless{}=255}-byte character-strings in RDATA order,
  \textbf{without separators}. Never concatenate across records.
\item
  Ignore \texttt{TXT} records that do not begin with \texttt{bitcoin:}.
\item
  If \textbf{more than one} \texttt{bitcoin:} record exists at the same
  name, the result is \textbf{invalid}.
\item
  The payload must parse as a valid BIP-321 URI and must contain
  \textbf{no private material}.
\item
  Do not cache longer than the DNS TTL.
\item
  Prefer displaying verified names as \texttt{₿user@domain}.
\end{itemize}

\hypertarget{dnssec-requirement-fail-closed}{%
\subsection{DNSSEC requirement (fail
closed)}\label{dnssec-requirement-fail-closed}}

BIP-353 payment instructions are only trustworthy if DNSSEC-validated.
SatsPath \textbf{does not trust an upstream resolver's AD bit} --- that
can be spoofed by a malicious resolver.

\texttt{DnssecPolicy} has two modes:

\begin{longtable}[]{@{}
  >{\raggedright\arraybackslash}p{(\columnwidth - 2\tabcolsep) * \real{0.3750}}
  >{\raggedright\arraybackslash}p{(\columnwidth - 2\tabcolsep) * \real{0.6250}}@{}}
\toprule\noalign{}
\begin{minipage}[b]{\linewidth}\raggedright
Mode
\end{minipage} & \begin{minipage}[b]{\linewidth}\raggedright
Behavior
\end{minipage} \\
\midrule\noalign{}
\endhead
\bottomrule\noalign{}
\endlastfoot
\texttt{Strict} (default) & Require real DNSSEC validation; \textbf{fail
closed} otherwise. \\
\texttt{DevInsecure} & Local testing only. Accept unvalidated records
with loud warnings. Never the default; requires
\texttt{-\/-allow-insecure-dns-for-dev}. \\
\end{longtable}

This build does not ship a local DNSSEC validator, so \texttt{Strict}
resolution \textbf{fails closed} with a clear message (option C of the
BIP-353 DNSSEC guidance). A future release can plug in a validating
resolver without changing the contract. The
\texttt{-\/-allow-insecure-dns-for-dev} flag prints a scary warning and
must never be used on mainnet.

\hypertarget{direct-txt-model-option-a}{%
\subsection{Direct TXT model (Option
A)}\label{direct-txt-model-option-a}}

The receiver controls their domain and publishes the record directly:

\begin{verbatim}
Name:
rodrigo.user._bitcoin-payment.satspath.dev

TXT:
bitcoin:?lno=lno1...
\end{verbatim}

\begin{itemize}
\tightlist
\item
  \textbf{Pros:} simple, direct, no platform dependency.
\item
  \textbf{Cons:} the receiver must update DNS whenever payment
  instructions change.
\end{itemize}

\hypertarget{managed-cnamedname-delegation-model-option-b}{%
\subsection{Managed CNAME/DNAME delegation model (Option
B)}\label{managed-cnamedname-delegation-model-option-b}}

The receiver delegates the BIP-353 label once to a SatsPath-managed
zone, which then maintains the \texttt{TXT}:

\begin{verbatim}
On example.com (added once by the receiver):
rodrigo.user._bitcoin-payment.example.com CNAME rodrigo.user._bitcoin-payment.satspath.dev

On satspath.dev (SatsPath-managed):
rodrigo.user._bitcoin-payment.satspath.dev TXT "bitcoin:?lno=lno1..."
\end{verbatim}

\begin{itemize}
\tightlist
\item
  Both the receiver's domain chain \textbf{and} the managed zone must be
  DNSSEC-signed.
\item
  The CNAME/DNAME chain must be DNSSEC-validated; resolution
  \textbf{fails closed} if not.
\item
  \textbf{Pros:} the receiver configures DNS once; SatsPath can rotate
  instructions safely.
\item
  \textbf{Cons:} requires trust in SatsPath managed-DNS availability;
  SatsPath must secure its DNS-provider API credentials (never
  committed).
\end{itemize}

\hypertarget{why-gmail-style-emails-cannot-be-verified-through-bip-353}{%
\subsection{Why Gmail-style emails cannot be verified through
BIP-353}\label{why-gmail-style-emails-cannot-be-verified-through-bip-353}}

BIP-353 requires publishing a DNS record under the \textbf{domain} of
the name. A user of \texttt{rodrigo@gmail.com} cannot publish
\texttt{rodrigo.user.\_bitcoin-payment.gmail.com} because they do not
control \texttt{gmail.com}.

\begin{quote}
SatsPath can resolve DNSSEC-backed BIP-353 payment instructions for
domains the receiver controls. For consumer email addresses like
\texttt{gmail.com}, SatsPath uses platform verification / the invite
flow instead, because the user cannot publish DNS records under
\texttt{gmail.com}.
\end{quote}

\hypertarget{rotating-payment-instructions}{%
\subsection{Rotating payment
instructions}\label{rotating-payment-instructions}}

DNS record changes require a \textbf{cryptographic identity-key
signature} --- never email login alone.

\begin{verbatim}
1. The receiver holds their SatsPath identity key locally.
2. The platform builds a challenge:
   satspath-dns-update:<alias>:<dns_name>:<sha256(new_uri)>:<nonce>:<expires_at>
3. The receiver signs the challenge with their identity key.
4. The platform verifies the signature against the profile's identity_pubkey.
5. The platform updates the DNS TXT record via the DnsPublisher adapter.
6. The platform records an audit entry:
   { alias, dns_name, old_uri_hash, new_uri_hash, timestamp, identity_fingerprint }.
\end{verbatim}

Email verification may gate \textbf{account access}, but DNS
payment-instruction changes always require identity proof.

\hypertarget{ttl-recommendations}{%
\subsection{TTL recommendations}\label{ttl-recommendations}}

\begin{longtable}[]{@{}ll@{}}
\toprule\noalign{}
Instruction kind & TTL \\
\midrule\noalign{}
\endhead
\bottomrule\noalign{}
\endlastfoot
Rotating on-chain address & 300 seconds \\
Reusable BOLT12 offer / Silent Payment & 1800 seconds \\
\end{longtable}

A direct on-chain address in a record is treated as \textbf{rotating}
and emits a privacy warning at publish time.

\hypertarget{bip-321-payload-handling}{%
\subsection{BIP-321 payload handling}\label{bip-321-payload-handling}}

A resolved record is parsed as BIP-321 payment instructions. Supported:

\begin{itemize}
\tightlist
\item
  on-chain address in the URI path
\item
  \texttt{amount} query parameter
\item
  \texttt{lightning} (BOLT11 invoice)
\item
  \texttt{lno} (BOLT12 offer)
\item
  \texttt{sp} (Silent Payment, \textbf{preview-only})
\item
  unknown non-\texttt{req-*} parameters are ignored safely
\item
  unknown \texttt{req-*} parameters make the URI \textbf{invalid}
\end{itemize}

\hypertarget{optional-satspath-profile-pointer-extension}{%
\subsubsection{Optional SatsPath profile pointer
extension}\label{optional-satspath-profile-pointer-extension}}

For richer SatsPath data without breaking BIP-321, a record may carry a
\textbf{non-required} key:

\begin{verbatim}
sp-profile=https%3A%2F%2Fsatspath.dev%2F.well-known%2Fsatspath%2Frodrigo
sp-profile-hash=<sha256>
\end{verbatim}

Wallets that don't understand SatsPath ignore it. SatsPath can use it to
fetch the signed profile \textbf{after} BIP-353 verifies the domain.

\hypertarget{security-checklist}{%
\subsection{Security checklist}\label{security-checklist}}

\begin{itemize}
\tightlist
\item[$\boxtimes$]
  DNSSEC required; \texttt{Strict} fails closed (no AD-bit trust).
\item[$\boxtimes$]
  Multiple \texttt{bitcoin:} records at one name → invalid.
\item[$\boxtimes$]
  Non-\texttt{bitcoin:} \texttt{TXT} records ignored.
\item[$\boxtimes$]
  Unknown \texttt{req-*} parameters → invalid.
\item[$\boxtimes$]
  Private material (\texttt{seed}, \texttt{xprv}/\texttt{tprv},
  \texttt{mnemonic}, \texttt{macaroon}, \texttt{cert},
  \texttt{api\_key}, \texttt{claim\_key}, \texttt{refund\_key},
  \texttt{preimage}, \ldots) rejected before publish and on resolve.
\item[$\boxtimes$]
  DNS updates require identity-key signatures; email-only rejected.
\item[$\boxtimes$]
  DNS-provider credentials never committed (trait +
  \texttt{MockDnsPublisher} only).
\item[$\boxtimes$]
  On-chain addresses use short TTL + rotation warning.
\item[$\boxtimes$]
  No funds moved, no signing, no broadcast --- \textbf{Mainnet Preview
  only}.
\end{itemize}

\hypertarget{example-mainnet-preview}{%
\subsection{Example: Mainnet Preview}\label{example-mainnet-preview}}

\begin{Shaded}
\begin{Highlighting}[]
\ExtensionTok{satspath}\NormalTok{ preview ₿rodrigo@satspath.dev 1000 }\AttributeTok{{-}{-}mainnet}
\end{Highlighting}
\end{Shaded}

Expected:

\begin{verbatim}
DNSSEC: valid
BIP-353: resolved
Payment URI: bitcoin:?lno=...
Mode: Mainnet Preview
No funds moved.
\end{verbatim}

\end{document}
